\chapter{Related Work}
\label{ch:related-work}
\raggedbottom

This chapter discusses approaches related to compiling coroutines and asynchronous code to WebAssembly,
as well as alternative designs for Stack-switching in Wasm.

\section{Approaches to Asynchronous Code in WebAssembly}

\subsection{Asyncify}

Before the Stack-switching proposal, the most widely used technique for compiling asynchronous languages to
WebAssembly was Asyncify~\cite{asyncify}, a binary transformation pass implemented in Binaryen.
Asyncify works by rewriting a compiled Wasm binary to make it possible to pause and resume execution at
arbitrary points. At suspension points, the transformation inserts code to unwind the current call stack
into a pre-allocated buffer, and at resumption points it inserts code to rewind the stack from that buffer.

The approach is analogous to the state-machine transformation used by the Kotlin compiler for the JVM and JS
targets: both convert stackful suspension into a representation that can be managed by a
runtime without native stack-switching support. The key difference is that Asyncify operates at the Wasm
binary level rather than at the source or IR level, making it language-agnostic.

Asyncify has been used to port a number of languages and runtimes to WebAssembly, including Python (Pyodide)
and Ruby. Its main drawbacks are increased binary size and runtime overhead from the stack save and restore
operations~\cite{asyncify}.

\subsection{JS Promise Integration (JSPI)}

An alternative approach for asynchronous WebAssembly is the JS Promise Integration (JSPI) proposal~\cite{jspi-proposal}.
JSPI provides a JavaScript API that allows Wasm imports and exports to be wrapped so that a
synchronous Wasm call can suspend waiting for a JavaScript \texttt{Promise} to settle, and then be resumed
when the promise resolves. From the Wasm side, JSPI is transparent: the module does not need to be modified.
The suspension and resumption are managed by the JavaScript engine.

JSPI addresses the specific use case of interoperating with asynchronous JavaScript APIs from synchronous
Wasm code. It does not provide general-purpose stack switching within the Wasm module itself, and therefore
cannot be used to implement language-level coroutines that suspend and resume within Wasm. The Stack-switching
proposal is a more general solution that covers this use case as well.

JSPI has been presented and discussed in the WebAssembly Stack Subgroup~\cite{jspi-wasm-meeting}. It has also
been partially implemented in SpiderMonkey using Stack-switching primitives~\cite{spidermonkey-stack-switching-diff}.

\subsection{WasmFX and Typed Continuations}

WasmFX~\cite{wasmfx-proposal} is an alternative proposal for adding stack-switching capabilities to
WebAssembly, based on the theory of algebraic effects and typed continuations.
Unlike the Stack-switching proposal (which uses a \texttt{cont.new}/\texttt{resume}/\texttt{suspend}
interface), WasmFX uses a \texttt{cont.new}/\texttt{cont.bind}/\texttt{resume}/\texttt{barrier}
interface with a different semantic model. The core difference is that WasmFX continuations are
designed to compose with effect handlers, enabling a richer set of control-flow abstractions.

WasmFX was presented and developed in the WebAssembly Stack Subgroup over several meetings~\cite{wasmfx-meeting-2021,wasmfx-meeting-2023}.
A prototype implementation in the Wasmtime engine was discussed in the subgroup in 2023~\cite{wasmfx-wasmtime-2023}.
The current Stack-switching proposal in the repository is a convergence of ideas from both proposals.

\section{Other Language Runtimes on WebAssembly}

Multiple language communities have explored how to compile their runtimes to WebAssembly, which requires
solving problems similar to the ones addressed in this thesis.

\textbf{Go and TinyGo.} The Go runtime uses goroutines, which are stackful coroutines. Two separate presentations
in the WebAssembly Stack Subgroup discussed the challenges of compiling Go to Wasm: one covering the full
Go runtime~\cite{go-wasm-meeting}, and one covering TinyGo~\cite{tinygo-wasm-meeting}, a Go compiler for
embedded and Wasm targets. Both noted the difficulty of mapping Go's concurrency model to WebAssembly's
linear stack model.

\textbf{OCaml.} Multicore OCaml adds support for algebraic effects, which generalize coroutines and other
control-flow abstractions. The OCaml community presented its experience compiling to Wasm in the
Stack Subgroup~\cite{ocaml-wasm-meeting}, discussing how native OCaml effects could be compiled to
WebAssembly using stack-switching primitives.

\textbf{Erlang.} The Erlang language model is based on lightweight processes, which are a form of
stackful coroutine. Challenges in implementing the Erlang runtime in WebAssembly were presented in the
Stack Subgroup~\cite{erlang-wasm-meeting}, noting that native Wasm stack switching would significantly
simplify the port compared to the Asyncify-based approach.

These cases, together with Kotlin, illustrate that language-level concurrency abstractions and
stack-switching support in WebAssembly are broadly needed across many language communities.

